\subsection*{Supplementary Note 2: additional results}
This note gives the detailed results that the main text summarises. Figure, table and Extended Data numbers refer to the main text.

\paragraph{Metric implementation}

After the probabilities of the released M2MRF checkpoints are stored as 16-bit maps, the authors' function returns \rn{e1.m2mrf.IDRiD.stored} on IDRiD and \rn{e1.m2mrf.DDR.stored} on DDR, between the value on the float output and the exact AUPR. The paper reports the mean of three runs (\rn{e1.m2mrf.IDRiD.paper} and \rn{e1.m2mrf.DDR.paper}). Applied to the probability maps of the retrained models,, the 11-threshold implementation changed \mAUPR{} on IDRiD by between $-$0.010 and +0.023 for the converged models, but on DDR, where the low-probability tail of the maps carries more weight, it added between 0.004 and 0.084 to them and 0.168 to the one model that had not converged: the inflation is model-dependent and is largest for the models whose maps are least calibrated, so it reorders models that differ by less than it (Fig.~\ref{M-fig:literature}d,e; per-model values in Supplementary Table~5).

The literature audit also recorded inconsistencies between copied tables and their sources, such as transposed lesion columns and a value carried from one metric into another; they are listed per paper in Supplementary Data~2. Of the method papers, \rn{lit.impl.exact} use the exact average precision, \rn{lit.impl.perimage} averages over images, and \rn{lit.pooled.stated} say whether pixels are pooled over the test set or scores are averaged over images. Only \rn{lit.rerun.all} papers state that every baseline was rerun, and \rn{lit.rerun.notstated} say nothing about where their comparison numbers come from. DeepLabv3+ is reported on IDRiD by \rn{lit.deeplab.n} papers with \mAUPR{} from \rn{lit.deeplab.lo} to \rn{lit.deeplab.hi}; copied values agree to the last digit, and the spread comes from the papers that reran a method under their own settings.

\paragraph{Author weights and author recipes}

Rerunning the authors' pipeline is the direct test of whether released code reproduces a reported number. With the authors' batch size of 4 and their 40,000 iterations, three seeds of the M2MRF rerun on IDRiD reached \rn{e1r.M2MRF.IDRiD.fn.lo}--\rn{e1r.M2MRF.IDRiD.fn.hi} with the authors' metric (mean \rn{e1r.M2MRF.IDRiD.fn}, s.d.\ \rn{e1r.M2MRF.IDRiD.fn.sd}) and \rn{e1r.M2MRF.IDRiD.exact.lo}--\rn{e1r.M2MRF.IDRiD.exact.hi} under our protocol (mean \rn{e1r.M2MRF.IDRiD.exact}), against 0.676 and 0.639 for the released checkpoint and 0.672 for the three-run mean printed in the paper. The models produced by the released training pipeline were therefore 5--7 points below the released weights on every seed, a difference about five times the seed spread of the rerun itself and of the same architecture under the unified protocol (s.d.\ 0.010); a first rerun with an effective batch of 1 in place of 4 is reported separately in Fig.~\ref{M-fig:original}c and Table~\ref{M-tab:original}.

The released HRDecoder checkpoints reached 0.707 on IDRiD and 0.461 on DDR under our protocol; they were trained by their authors on a canvas twice as large as the unified one and on three to five times as many pixels. Rerun once with the authors' recipe (their geometry, augmentation, schedule and all 54 training images; one process in place of four), HRDecoder reached \rn{e1r.HRDecoder.IDRiD.fn} with the authors' metric and \rn{e1r.HRDecoder.IDRiD.exact} under our protocol on IDRiD, 0.04 and 0.03 below the released checkpoint (0.713 and 0.707); the checkpoint selection of the authors' pipeline, which evaluates on the test split, would have added 0.002 to the rerun and does not account for the difference. Applied to the independent test sets as released, the four checkpoints behaved like the retrained models (Fig.~\ref{M-fig:original}d): the DDR-trained checkpoints scored higher on the other pixel-labelled standard-field sources (IDRiD, the standard-field TJDR set) than on their own test set, kept 74--83\% of their internal score on MAPLES-DR, and kept 36--60\% of it on the coarsely or regionally labelled sets and on the ultra-widefield images, whereas the IDRiD-trained checkpoints fell on every other source (to 0.06--0.48 from 0.64 and 0.71).

\paragraph{Internal tests}

On the internal tests (Fig.~\ref{M-fig:retrained}a and Table~\ref{M-tab:main}), adapted foundation models and a general-purpose high-resolution network led both sources, and the differences among them were within the seed-to-seed variation. HRNet-W48 reached 0.700 $\pm$ 0.004 \mAUPR{} on IDRiD and 0.447 $\pm$ 0.009 on DDR (mean $\pm$ s.d.\ over three seeds), SAMed, a low-rank adaptation of SAM ViT-B trained on 512-pixel windows of the 1,440-pixel canvas, 0.706 $\pm$ 0.015 and 0.460 $\pm$ 0.027, and RETFound with a UPerNet head 0.699 $\pm$ 0.009 and 0.450 $\pm$ 0.018; the seed-0 margins of the two foundation models over HRNet-W48 on DDR (0.026 and 0.024) shrank to 0.013 and 0.004 over three seeds, within the seed scatter of the models themselves, and on IDRiD the three-seed margins (0.006 and $-$0.001) are within the width of the 27-image test set. SAM-Adapter, run with one seed, was level with SAMed on IDRiD (0.706) and 0.03 below it on DDR (0.433), and Swin-Unet (0.688 $\pm$ 0.003 and 0.431 $\pm$ 0.001) and the self-configuring nnU-Net (0.665 $\pm$ 0.008 and 0.431 $\pm$ 0.003) followed within 0.04 on both sources. The seven fundus-specific networks were below them: CLC-Net, the highest of the family on DDR (0.661 and 0.414), was 0.02 above U-Net on IDRiD and 0.04 above it on DDR, M2MRF reached 0.657 $\pm$ 0.010 and 0.369 $\pm$ 0.011, within 0.02 of its released checkpoint on IDRiD (0.639) and below it on DDR (0.401), and the other five spanned 0.54--0.65 on IDRiD and 0.23--0.36 on DDR (Table~\ref{M-tab:main}). The difference between the first two models of each column, less than 0.001 on IDRiD and \rn{bstd.DDR.gap12} on DDR, is at most of the order of the median seed standard deviation of the table (\rn{bstd.IDRiD.sd.median} and \rn{bstd.DDR.sd.median}), whereas the span of the table (0.41 on IDRiD, 0.28 on DDR) is thirty to forty times larger. The 95\% bootstrap intervals over test images (seed 0; Supplementary Table~4) are wider still: their median width across models is \rn{ci.IDRiD.width.lo}--\rn{ci.IDRiD.width.hi} AUPR for microaneurysms, hemorrhages and hard exudates on IDRiD and \rn{ci.DDR.width.lo}--\rn{ci.DDR.width.hi} on DDR, and \rn{ci.IDRiD.SE.width} and \rn{ci.DDR.SE.width} for soft exudates, which appear in few test images; the interval of the four-class mean, from the same resamples, has a median width of \rn{mci.IDRiD.width.median} on IDRiD and \rn{mci.DDR.width.median} on DDR, so the intervals of the leading models overlap (SAMed \rn{mci.samed.DDR}, RETFound with UPerNet \rn{mci.retfound.DDR}, HRNet-W48 \rn{mci.hrnet.DDR} on DDR). IDRiD has 27 test images, and its intervals are correspondingly wide; we use it to connect to the literature and base comparisons between architectures on DDR and on the independent tests.

Lesion classes differed more than models. The median AUPR over the \rn{n.main} models was 0.87 for hard exudates and 0.44 for microaneurysms on IDRiD, and 0.60 and 0.15 on DDR; on DDR the highest microaneurysm score was 0.189 (0.14--0.24) for RETFound with UPerNet (seed 0; 95\% bootstrap interval in brackets), against 0.155 (0.12--0.20) for HRNet-W48 and 0.154 (0.11--0.17) for nnU-Net, and the highest soft-exudate score 0.550 (0.42--0.65) for HRNet-W48; the intervals of the leading models overlap on every class.

\paragraph{Case study: one backbone, three decoders}

HRDecoder, M2MRF and HRNet-W48 share the HRNet-W48 backbone, so their three rows compare a general head with two lesion-specific decoders on one backbone. Under the unified protocol the general head led on both sources (0.700 and 0.447), M2MRF followed (0.657 and 0.369) and HRDecoder was third (0.647 and 0.296), a span of 0.05 on IDRiD and 0.15 on DDR, whereas the released HRDecoder checkpoints score 0.707 and 0.461 under the same evaluation (see above). Three further runs separate the explanations. The authors' recipe, rerun once, gave 0.674 on IDRiD under our protocol, 0.03 above the unified run and 0.03 below the released checkpoint. Doubling the training amount of the unified protocol at the 1,440-pixel canvas did not close the gap (0.615 on IDRiD against 0.647 for the single budget, and 0.304 on DDR against 0.296, one seed each), whereas the same doubling raised the shared-backbone control HRNet-W48 on DDR from 0.447 to 0.469. The unified protocol at the authors' 2,880-pixel canvas and 1,920-pixel window, with the training amount counted in equivalent epochs at that canvas, gave \rn{e7.hrdecoder.2880.aupr}, within 0.01 of the 1,440-pixel run (\rn{e7.hrdecoder.1440.aupr}); hard exudates and microaneurysms gained (\rn{e7.hrdecoder.2880.EX} against \rn{e7.hrdecoder.1440.EX}, \rn{e7.hrdecoder.2880.MA} against \rn{e7.hrdecoder.1440.MA}) and hemorrhages and soft exudates lost (\rn{e7.hrdecoder.2880.HE} against \rn{e7.hrdecoder.1440.HE}, \rn{e7.hrdecoder.2880.SE} against \rn{e7.hrdecoder.1440.SE}). Neither the canvas nor the training amount therefore accounts for the released checkpoint; what the authors' recipe adds beyond them, 0.03--0.06 on one seed, lies in the parts of the recipe that the protocol replaces (rotations of $\pm45^\circ$ without photometric distortion, training on all 54 images without a validation split, and the last-step checkpoint; Supplementary Table~3).

\paragraph{Foundation models}

Zero-shot use did not segment DR lesions. The automatic masks of SAM and SAM~2, with the density and area cap chosen on the validation pool, reached an any-lesion AUPR of \rn{fm.any.sam.IDRiD} and \rn{fm.any.samtwo.IDRiD} on IDRiD and \rn{fm.any.sam.DDR} on DDR, against \rn{fm.any.ref.nnunet.IDRiD} for nnU-Net's maps on the same metric, and even when each mask was given the class it overlaps most in the ground truth the \mAUPR{} stayed at \rn{fm.zero.ub.value}. Text prompts to SAM~3 found two of the four classes: with the wording chosen on the validation pool, hard exudates reached an AUPR of \rn{fm.sam3.EX.IDRiD} on IDRiD and hemorrhages \rn{fm.sam3.HE.IDRiD} (0.24 and 0.23 on DDR), soft exudates \rn{fm.sam3.SE.IDRiD} on IDRiD and 0.02 on DDR, and microaneurysms \rn{fm.sam3.MA.IDRiD} under every wording. The wording mattered as much as the class, from \rn{fm.sam3.EX.min} to \rn{fm.sam3.EX.max} for hard exudates and from \rn{fm.sam3.SE.min} to \rn{fm.sam3.SE.max} for soft exudates, where the clinical term (``cotton wool spot'') scored lowest and a descriptive wording (``white fluffy patch'') highest (Fig.~\ref{M-fig:foundation}c). Prompts derived from the ground truth tell where the gap lies. With a box around every ground-truth component, SAM~2 reached \rn{fm.oracle.all} \mAUPR{} on IDRiD, but an IDRiD image carries on average 44 microaneurysm and 142 hard-exudate components, and with 1, 5 and 20 boxes per image and class the \mAUPR{} was \rn{fm.oracle.one}, \rn{fm.oracle.five} and \rn{fm.oracle.twenty} (Fig.~\ref{M-fig:foundation}b); one positive point per component gave \rn{fm.point.all.IDRiD} on IDRiD and \rn{fm.point.all.DDR} on DDR. Given every box, the supervised models' own maps cut to the union of the same boxes scored \rn{fm.boxref.nnunet.IDRiD} (nnU-Net) and \rn{fm.boxref.m2mrf.IDRiD} (M2MRF) on IDRiD and \rn{fm.boxref.nnunet.DDR} and \rn{fm.boxref.m2mrf.DDR} on DDR, against \rn{fm.box.sam.IDRiD} and \rn{fm.box.sam.DDR} for SAM ViT-B, \rn{fm.box.samtwo.IDRiD} and \rn{fm.box.samtwo.DDR} for SAM~2 and \rn{fm.box.medsam.IDRiD} and \rn{fm.box.medsam.DDR} for MedSAM: inside a box the general promptable models delineated the lesion as well as the best supervised model on IDRiD and better on DDR. The gap of the promptable models therefore lies in finding lesions, not in drawing them.

The pretraining of the encoder was isolated under one UPerNet head, one recipe and one training amount, in two pairs fixed in advance: an encoder pretrained on retinal images against its own starting point, under each of two pretraining objectives (Fig.~\ref{M-fig:foundation}d,e; Supplementary Table~10). On DDR, RETFound's masked-autoencoder encoder reached \rn{fm.enc.rm.int} internally and \rn{fm.enc.rm.TJDR-std} on the standard-field TJDR set, below the ImageNet masked autoencoder it was initialised from (\rn{fm.enc.mi.int} and \rn{fm.enc.mi.TJDR-std}), and RETFound's DINOv2 encoder \rn{fm.enc.rd.int} and \rn{fm.enc.rd.TJDR-std} against \rn{fm.enc.di.int} and \rn{fm.enc.di.TJDR-std} for DINOv2 itself; the two pairs differed by $-$\rn{fm.enc.c1.int.abs} and $-$\rn{fm.enc.c2.int.abs} internally and the second pair by less than 0.015 on each of the three fine-grained independent sets. On IDRiD the same two pairs gave \rn{fm.enc.i.rm.int} against \rn{fm.enc.i.mi.int} and \rn{fm.enc.i.rd.int} against \rn{fm.enc.i.di.int} ($-$\rn{fm.enc.i.c1.int.abs} and $-$\rn{fm.enc.i.c2.int.abs}). Under one head, one recipe and one seed per encoder, retinal pretraining therefore added nothing to either starting point under either objective. The two natural-image encoders, ViT-L models that have never seen a fundus photograph, exceeded every row of Table~\ref{M-tab:main} on both sources (\rn{fm.enc.mi.int} and \rn{fm.enc.di.int} on DDR against 0.469 for the best ranked model on seed 0; \rn{fm.enc.i.mi.int} and \rn{fm.enc.i.di.int} on IDRiD against 0.717), by a margin larger than the seed scatter of any ranked model. These encoder cells were run as controls with one seed and one learning rate each and are reported in Supplementary Table~10 rather than as ranked rows.

\paragraph{Independent tests}

Moving to a source that was never used for training did not by itself cost performance; the annotation tier and the imaging field did (Fig.~\ref{M-fig:retrained}b,c and Table~\ref{M-tab:independent}; the \rn{n.main} models on all eight test sets, seed means as in Table~\ref{M-tab:main}). The models trained on DDR averaged \rn{indep.DDR.int.mean} \mAUPR{} on the DDR test set and scored higher on the two other standard-field sources with pixel-level labels, \rn{indep.DDR.IDRiD.mean} on IDRiD and \rn{indep.DDR.TJDR-std.mean} on the standard-field TJDR set (every model higher; two-sided Wilcoxon signed-rank $P$ \rn{indep.DDR.TJDR-std.p}), and the same on MAPLES-DR (\rn{indep.DDR.MAPLES-DR.mean}; \rn{indep.DDR.MAPLES-DR.nlower} of \rn{n.main} lower, $P$ \rn{indep.DDR.MAPLES-DR.p}). They lost a third to a half of their internal score wherever the labels or the geometry differ from the training source: \rn{indep.DDR.FGADR.mean} on FGADR (\rn{indep.DDR.FGADR.pct}\%) and \rn{indep.DDR.Retinal-Lesions.mean} on Retinal-Lesions (\rn{indep.DDR.Retinal-Lesions.pct}\%), whose lesions are drawn as coarse regions, \rn{indep.DDR.DiaRetDB1.mean} pixel-level \mAUPR{} against the region-level consensus of DiaRetDB1 (\rn{indep.DDR.DiaRetDB1.pct}\%), and \rn{indep.DDR.TJDR-uwf.mean} on the 133$^\circ$ ultra-widefield TJDR images (\rn{indep.DDR.TJDR-uwf.pct}\%; every model lower on each of the four sets, $P$ \rn{indep.DDR.TJDR-uwf.p}). The models trained on the 44 IDRiD images started higher internally (mean \rn{indep.IDRiD.int.mean}) and ended lower on every other source, fine-grained or not: \rn{indep.IDRiD.DDR.mean} on DDR, \rn{indep.IDRiD.TJDR-std.mean} on the standard-field TJDR set and \rn{indep.IDRiD.MAPLES-DR.mean} on MAPLES-DR (\rn{indep.IDRiD.TJDR-std.pct}\% to \rn{indep.IDRiD.DDR.pct}\%), and \rn{indep.IDRiD.Retinal-Lesions.mean}--\rn{indep.IDRiD.DiaRetDB1.mean} on the coarse, region-level and ultra-widefield sets (\rn{indep.IDRiD.DiaRetDB1.pct}\% to \rn{indep.IDRiD.Retinal-Lesions.pct}\%; every model lower on every set, $P$ \rn{indep.IDRiD.FGADR.p}). The 0.26 advantage of IDRiD over DDR on the internal tests therefore reversed outside the training source: on the six independent sets the DDR-trained checkpoint of a model was better than its IDRiD-trained one in \rn{indep.pairs.ddrbetter} of \rn{indep.pairs.n} pairs, in \rn{indep.pairs.std.lo}--\rn{indep.pairs.std.hi} of \rn{n.main} on each of the five standard-field sets; only on the ultra-widefield images were the IDRiD-trained checkpoints usually ahead (\rn{indep.pairs.TJDR-uwf.ddrbetter} of \rn{n.main} DDR-trained ahead), consistent with the larger field and higher resolution of their training images. For both training sources the pixel-labelled standard-field sets were the least affected (a gain from DDR, the smallest loss from IDRiD) and the coarse, region-labelled and ultra-widefield sets the most. On the ultra-widefield images the inference scale mattered less than the field itself: resizing them to the 1,440-pixel canvas instead of predicting at native resolution with sliding windows, the comparison variant of the registered plan, moved the median \mAUPR{} by +0.02 from both sources (13 and 14 of \rn{n.ranked} models higher at the smaller scale, differences from $-$0.12 to +0.13), far from the 0.2--0.4 that separate this set from the standard-field ones (Extended Data Fig.~\ref{M-edfig:uwf}b).

Internal rank predicted independent rank where the labels were drawn the same way. For DDR-trained models the Spearman correlation between internal \mAUPR{} and the score on IDRiD, the standard-field TJDR set and MAPLES-DR was \rn{indep.DDR.IDRiD.rho}, \rn{indep.DDR.TJDR-std.rho} and \rn{indep.DDR.MAPLES-DR.rho}: HRNet-W48, third internally, was first on IDRiD (\rn{indep.DDR.IDRiD.first.value}) and on the standard-field TJDR set (\rn{indep.DDR.TJDR-std.first.value}), SAMed, first internally, on MAPLES-DR (\rn{indep.DDR.MAPLES-DR.first.value}) and FGADR (\rn{indep.DDR.FGADR.first.value}; $\rho$ = \rn{indep.DDR.FGADR.rho}), and SAMed also led the mean of the three fine-grained independent sets (\rn{indep.DDR.fine.first.value}; HRNet-W48 \rn{indep.DDR.fine.hrnet}). On Retinal-Lesions and DiaRetDB1 the correlation was \rn{indep.DDR.Retinal-Lesions.rho} and \rn{indep.DDR.DiaRetDB1.rho} but the first place went to DeepLabv3 (\rn{indep.DDR.Retinal-Lesions.first.value} and \rn{indep.DDR.DiaRetDB1.first.value}), tenth internally, whose output stride of 8 yields region-sized predictions, and on the ultra-widefield images the correlation vanished ($\rho$ = \rn{indep.DDR.TJDR-uwf.rho}) and the fully fine-tuned MedSAM, ninth internally, was first (\rn{indep.DDR.TJDR-uwf.first.value}). For IDRiD-trained models the correlation was \rn{indep.IDRiD.DDR.rho} on DDR and \rn{indep.IDRiD.MAPLES-DR.rho} on MAPLES-DR and \rn{indep.IDRiD.TJDR-uwf.rho}--\rn{indep.IDRiD.FGADR.rho} elsewhere; SAMed led DDR, the standard-field TJDR set and FGADR (\rn{indep.IDRiD.DDR.first.value}, \rn{indep.IDRiD.TJDR-std.first.value}, \rn{indep.IDRiD.FGADR.first.value}) and the fine-grained mean (\rn{indep.IDRiD.fine.first.value}), whereas TC-Net, tenth internally, led MAPLES-DR (\rn{indep.IDRiD.MAPLES-DR.first.value}), DeepLabv3, eighteenth, led Retinal-Lesions and DiaRetDB1 (\rn{indep.IDRiD.Retinal-Lesions.first.value} and \rn{indep.IDRiD.DiaRetDB1.first.value}), and MedSAM, fourteenth, the ultra-widefield images (\rn{indep.IDRiD.TJDR-uwf.first.value}). RETFound with UPerNet, second internally on DDR, scored lowest among the converged models on the ultra-widefield images from DDR (0.111): its 224-pixel windows see 1/19 of the 4,256-pixel canvas at a time.

Internal and independent performance were correlated across models (Spearman $\rho$ = \rn{indep.DDR.rho} for DDR-trained and \rn{indep.IDRiD.rho} for IDRiD-trained models, internal \mAUPR{} against the mean of the three fine-grained independent sets), and the model that led the fine-grained independent sets, \rn{indep.DDR.top.model} (\rn{indep.DDR.top.value} from DDR, \rn{indep.IDRiD.top.value} from IDRiD), was \rn{indep.DDR.sametop} as on the internal test. Averaged by family over both sources, the adapted foundation models were ahead on internal (\rn{family.foundation.internal}) and on fine-grained independent tests (\rn{family.foundation.independent}), the self-configuring baseline followed (\rn{family.auto.internal} and \rn{family.auto.independent}), and the fundus-specific (\rn{family.fundus.internal} and \rn{family.fundus.independent}) and general networks (\rn{family.general.internal} and \rn{family.general.independent}) were level with each other (Fig.~\ref{M-fig:retrained}d).

\paragraph{Training source and annotation granularity}

The sources differ in how lesions were drawn before they differ in anything a model does (Fig.~\ref{M-fig:sources}a; Methods). Counted as connected components of the prepared masks over every image of each source, an IDRiD image carries on average 43 microaneurysm and 144 hard-exudate components whose median diameter is 0.5\% of the field-of-view width, and a DDR image 14 and 31 at 0.4\%; MAPLES-DR, drawn pixel by pixel by retinologists, has 27 microaneurysms and 16 hard exudates per image at 0.8\% and 1.0\%. The two TJDR sets carry few components per image (3--7 per class) of the same small size (0.4--1.8\%), so their difference to IDRiD and DDR is in how many lesions the images carry, not in how they were outlined. In the coarse and region-consensus sources the components are few and large: a Retinal-Lesions image carries 2 hard-exudate and 8 hemorrhage components with median diameters of 4.2\% and 3.5\% of the field, and a DiaRetDB1 image 1.5 and 4.5 at 4.1\% and 2.9\%, eight to nine times the hard-exudate diameter of IDRiD and DDR and two to four times their hemorrhage diameter, which for the same retina means that several lesions and the tissue between them were enclosed by one outline. FGADR, drawn as a mixture of single lesions and regions, sits between: 29 hard-exudate and 21 hemorrhage components per image, as many as DDR's, at 0.8\% and 1.2\%, somewhat larger than DDR's and a fifth to a third of Retinal-Lesions'.

To separate the effect of the training source from that of the model, the five representative models (U-Net, nnU-Net, M2MRF, SAM2-UNet and RETFound with UPerNet) were also trained on the standard-field TJDR images and on FGADR, at the pixel budget of DDR and with one seed, which together with the main table gives a matrix of four training sources by eight test sets (Fig.~\ref{M-fig:sources}b,d; for the TJDR- and FGADR-trained models, whose sources fix no single field of view, the ultra-widefield set is predicted at the field-of-view scale, Methods). Every source served its own test set best: in-source training gave \rn{src.idrid.idrid} \mAUPR{} on IDRiD, \rn{src.ddr.ddr} on DDR, \rn{src.tjdr.tjdr} on the standard-field TJDR set and \rn{src.fgadr.fgadr} on FGADR, against \rn{src.idrid.foreign.lo}--\rn{src.idrid.foreign.hi}, \rn{src.ddr.foreign.lo}--\rn{src.ddr.foreign.hi}, \rn{src.tjdr.foreign.lo}--\rn{src.tjdr.foreign.hi} and \rn{src.fgadr.foreign.lo}--\rn{src.fgadr.foreign.hi} from the other three sources (Fig.~\ref{M-fig:sources}c). Training elsewhere cost \rn{cross.mean} \mAUPR{} on average over the twelve off-diagonal cells and up to \rn{cross.max} (\rn{cross.max.train}-trained models on \rn{cross.max.test}). Which foreign source was best depended on the labels of the test set rather than on the number of training images: DDR, the largest fine-grained source (383 images), was the best foreign source on every standard-field pixel-level test set (\rn{src.ddr.idrid} on IDRiD, \rn{src.ddr.tjdr} on the standard-field TJDR set, \rn{src.ddr.maples} on MAPLES-DR), whereas FGADR, three times larger (1,096 images) but drawn as mixed regions, was the best foreign source only on the two sets with coarse or region labels (\rn{src.fgadr.retles} on Retinal-Lesions and \rn{src.fgadr.diaret} on DiaRetDB1, against \rn{src.ddr.retles} and \rn{src.ddr.diaret} from DDR) and on the ultra-widefield images (\rn{src.fgadr.uwf} against \rn{src.ddr.uwf} from DDR; its cells are predicted at the field-of-view scale, Methods), and lost \rn{cross.fgadr.fine} on average on the three fine-grained sources relative to in-source training; the standard-field TJDR set, the smallest source (172 images), transferred worst (\rn{src.tjdr.idrid} on IDRiD, \rn{src.tjdr.ddr} on DDR).

We then added mixed or coarse labels to fine-label training while keeping the number of patch pixels fixed, for U-Net and M2MRF with one seed; differences smaller than twice the median seed standard deviation of the main table (0.018 on DDR) are reported without interpretation (Fig.~\ref{M-fig:sources}e and Table~\ref{M-tab:factors}). On the DDR test set, adding the 1,096 FGADR images to DDR changed \mAUPR{} by \rn{e4.unet.fgadr.ddr.delta} for U-Net (\rn{e4.unet.ref.ddr} to \rn{e4.unet.fgadr.ddr}) and \rn{e4.m2mrf.fgadr.ddr.delta} for M2MRF (\rn{e4.m2mrf.ref.ddr} to \rn{e4.m2mrf.fgadr.ddr}); adding only 383 FGADR images, as many as DDR has, by \rn{e4.unet.fgadr383.ddr.delta} and \rn{e4.m2mrf.fgadr383.ddr.delta} (\rn{e4.unet.fgadr383.ddr} and \rn{e4.m2mrf.fgadr383.ddr}), so the loss is not a matter of the fine images being outnumbered; adding Retinal-Lesions by \rn{e4.unet.retles.ddr.delta} and \rn{e4.m2mrf.retles.ddr.delta} (\rn{e4.unet.retles.ddr} and \rn{e4.m2mrf.retles.ddr}). On the IDRiD test set the same three training sets moved U-Net by \rn{e4.unet.fgadr.idrid.delta}, \rn{e4.unet.fgadr383.idrid.delta} and \rn{e4.unet.retles.idrid.delta} (from \rn{e4.unet.ref.idrid}) and M2MRF by \rn{e4.m2mrf.fgadr.idrid.delta}, \rn{e4.m2mrf.fgadr383.idrid.delta} and \rn{e4.m2mrf.retles.idrid.delta} (from \rn{e4.m2mrf.ref.idrid}), and on the standard-field TJDR and MAPLES-DR sets U-Net by \rn{e4.unet.other.lo} to \rn{e4.unet.other.hi} and M2MRF by \rn{e4.m2mrf.other.lo} to \rn{e4.m2mrf.other.hi}. Mixed labels added to fine ones were therefore at best neutral and at worst costly, and coarse labels cost M2MRF \rn{e4.m2mrf.retles.cost.lo}--\rn{e4.m2mrf.retles.cost.hi} on every standard-field fine-grained test set. Trained on Retinal-Lesions alone, the two models scored \rn{e4.unet.retlesonly.retles} and \rn{e4.m2mrf.retlesonly.retles} on its own test split but \rn{e4.unet.retlesonly.ddr} and \rn{e4.m2mrf.retlesonly.ddr} on DDR and \rn{e4.unet.retlesonly.idrid} and \rn{e4.m2mrf.retlesonly.idrid} on IDRiD; trained on FGADR alone, \rn{e4.unet.fgadronly.idrid} and \rn{e4.m2mrf.fgadronly.idrid} on IDRiD.

A revision of the labels themselves, as opposed to a change of source, moved the scores far less. A refined annotation of the 27 IDRiD test images was released in 2026 (Refined IDRiD; one class per pixel, 1,024-pixel crops of the retinal region). We located each crop in the original image (every crop is the retinal bounding box padded to a square and resized; normalised cross-correlation with the original $\geq$ 0.9988 for all 27), resampled the official labels through the same geometry, and scored the test maps of the \rn{ref.n} IDRiD-trained models of Table~\ref{M-tab:main} against both label sets without retraining or re-inference. The revision changed 13 of the 27 images for microaneurysms and hemorrhages, 14 for hard exudates and 6 of 14 for soft exudates, moving 5.7\%, 1.6\%, 4.0\% and 0.7\% of the labelled pixels (pooled intersection over union 0.94, 0.98, 0.96 and 0.99) and no component count except hard exudates (3,895 to 3,882). Under the same protocol the \mAUPR{} of every model changed by between \rn{ref.delta.min} and \rn{ref.delta.max} (microaneurysm AUPR by \rn{ref.MA.delta.mean} on average, the other classes by less than 0.002), and the order of the \rn{ref.n} models was unchanged (Kendall $\tau$ = \rn{ref.tau}). The 1,024-pixel re-scoring itself stayed within \rn{ref.geometry.absmax} \mAUPR{} of the main-table test at the 1,440-pixel canvas ($\tau$ = \rn{ref.geometry.tau}). What separates the sources in Fig.~\ref{M-fig:sources} is therefore how lesions were defined and outlined, which a correction of individual masks within one style does not touch (Supplementary Table~9).

\paragraph{Lesion size, input resolution and training amount}

Lesion size is the property that separates the four classes most clearly (Fig.~\ref{M-fig:scale}a). On the training images the median diameter of a microaneurysm is 0.5\% of the field-of-view diameter and that of a hemorrhage 1.1\%; we froze these two values, computed on training images only, as the boundaries between small, medium and large lesions before the main-table training. Resizing a whole image to a small input removes lesions from the training target itself: at 224 pixels \rn{vanish.IDRiD.224.MA}\% of the microaneurysms and \rn{vanish.IDRiD.224.EX}\% of the hard exudates of the IDRiD training set no longer cover a single pixel, and on DDR \rn{vanish.DDR.224.MA}\% and \rn{vanish.DDR.224.EX}\%; at 384 pixels the DDR figures are \rn{vanish.DDR.384.MA}\% and \rn{vanish.DDR.384.EX}\%, and on the 1,440-pixel canvas of the unified protocol essentially none are lost (Fig.~\ref{M-fig:scale}b).

Lesion-level recall followed size for every model (Fig.~\ref{M-fig:scale}c). On the DDR test set the median recall over the \rn{n.ranked} models (seed 0, threshold chosen on the validation set) was 0.26 for small, 0.49 for medium and 0.66 for large lesions; the best models detected 0.30 of the small ones (Swin-Unet and SAM2-UNet, nnU-Net ResEnc-M 0.29) and DeepLabv3, decoding at an output stride of 8, 0.07. By class, the median recall of small lesions was 0.46 for microaneurysms, 0.34 for hard exudates and 0.15 for hemorrhages, and the median lesion-level precision 0.29, 0.73 and 0.65. Input resolution acted mainly through this channel. Shrinking the canvas from 1,440 to 512 pixels at an equal pixel budget lowered \mAUPR{} on DDR from \rn{e7.m2mrf.1440.aupr} to \rn{e7.m2mrf.512.aupr} for M2MRF and from \rn{e7.retfound.1440.aupr} to \rn{e7.retfound.512.aupr} for RETFound with UPerNet (one seed), the recall of small lesions from \rn{e7.m2mrf.1440.small} to \rn{e7.m2mrf.512.small} and from \rn{e7.retfound.1440.small} to \rn{e7.retfound.512.small}, and the microaneurysm AUPR from \rn{e7.m2mrf.1440.MA} to \rn{e7.m2mrf.512.MA} and from \rn{e7.retfound.1440.MA} to \rn{e7.retfound.512.MA} (Fig.~\ref{M-fig:scale}d). Patches at native resolution, where a 960-pixel patch covers a smaller part of the retina and the fixed pixel budget buys fewer passes over each image, raised the small-lesion recall of M2MRF to \rn{e7.m2mrf.native.small} and of RETFound to \rn{e7.retfound.native.small}, but lowered \mAUPR{} to \rn{e7.m2mrf.native.aupr} and \rn{e7.retfound.native.aupr}: for RETFound the microaneurysm and hemorrhage AUPR rose (\rn{e7.retfound.1440.MA} to \rn{e7.retfound.native.MA} and \rn{e7.retfound.1440.HE} to \rn{e7.retfound.native.HE}) while the hard- and soft-exudate AUPR fell (\rn{e7.retfound.1440.EX} to \rn{e7.retfound.native.EX} and \rn{e7.retfound.1440.SE} to \rn{e7.retfound.native.SE}), so the two smallest classes gained from native resolution and the two largest lost.

More training images helped, and the two models gained at different points of the curve. With 10\%, 50\% and 100\% of the DDR training set (38, 192 and 383 images; equal number of patch pixels, so each image is seen ten, two and one times as often), \mAUPR{} on the DDR test set was \rn{e8.unet.10.ddr}, \rn{e8.unet.50.ddr} and \rn{e8.unet.100.ddr} for U-Net and \rn{e8.m2mrf.10.ddr}, \rn{e8.m2mrf.50.ddr} and \rn{e8.m2mrf.100.ddr} for M2MRF, and on IDRiD \rn{e8.unet.10.idrid}, \rn{e8.unet.50.idrid} and \rn{e8.unet.ref.idrid} and \rn{e8.m2mrf.10.idrid}, \rn{e8.m2mrf.50.idrid} and \rn{e8.m2mrf.ref.idrid} (Fig.~\ref{M-fig:scale}e): U-Net had most of its gain with the first half of DDR, M2MRF with the second half. Merging the three fine-grained sources IDRiD, DDR and standard-field TJDR (599 images) at the same budget gave \rn{e8.unet.fine.ddr} and \rn{e8.m2mrf.fine.ddr} on DDR, within 0.02 of DDR alone, and raised the scores on the sets outside the training sources by 0.00--0.04 (MAPLES-DR \rn{e8.unet.fine.maples} and \rn{e8.m2mrf.fine.maples} against \rn{e8.unet.ref.maples} and \rn{e8.m2mrf.ref.maples}; FGADR \rn{e8.unet.fine.fgadr} and \rn{e8.m2mrf.fine.fgadr} against \rn{e8.unet.ref.fgadr} and \rn{e8.m2mrf.ref.fgadr}; Retinal-Lesions \rn{e8.unet.fine.retles} and \rn{e8.m2mrf.fine.retles} against \rn{e8.unet.ref.retles} and \rn{e8.m2mrf.ref.retles}), the gain of more fine-grained images rather than of in-source data. A source merged into the training set raised the score on its own test set by more, 0.05--0.07 (standard-field TJDR \rn{e8.unet.fine.tjdr} and \rn{e8.m2mrf.fine.tjdr} against \rn{e8.unet.ref.tjdr} and \rn{e8.m2mrf.ref.tjdr}; IDRiD \rn{e8.unet.fine.idrid} and \rn{e8.m2mrf.fine.idrid} against \rn{e8.unet.ref.idrid} and \rn{e8.m2mrf.ref.idrid}), which is the in-source gain of the training-source matrix and not generalization; these two tests are in-source for the merged run and are marked as such in Fig.~\ref{M-fig:scale}e. Exchanging the loss of U-Net with the architecture fixed moved \mAUPR{} on DDR between 0.357 (cross-entropy alone) and 0.378 (cross-entropy with batch Dice, the main-table recipe; seed 0, as for every loss cell); Dice alone gave 0.237, the per-image Dice of the registered recipe 0.119, the rare-class collapse described in the main text (Extended Data Fig.~\ref{M-edfig:loss}).

\paragraph{Disease stage of the test images}

Scores also depend on the disease stage of the test images, because the pooled AUPR follows the number of lesion pixels available. On the four test sets that carry a DR grade per image, every score was recomputed within grades 0--1 and within grades 2--4 with the thresholds of the full test (Methods; Supplementary Table~12). On the DDR test set the median \mAUPR{} of the DDR-trained models (seed 0) was \rn{clin.ddr.lo.maupr} on the \rn{clin.ddr.lo.n} images of grades 0--1 and \rn{clin.ddr.hi.maupr} on the \rn{clin.ddr.hi.n} images of grades 2--4, on FGADR \rn{clin.fgadr.lo.maupr} and \rn{clin.fgadr.hi.maupr} (\rn{clin.fgadr.lo.n} and \rn{clin.fgadr.hi.n} images) and on Retinal-Lesions \rn{clin.retles.lo.maupr} and \rn{clin.retles.hi.maupr} (\rn{clin.retles.lo.n} and \rn{clin.retles.hi.n}); the image-level AUC on DDR was \rn{clin.ddr.lo.auc} within grades 0--1, where few images carry a lesion of the class (defined for \rn{clin.ddr.lo.auc.n} of the \rn{n.main} models), and \rn{clin.ddr.hi.auc} within grades 2--4. MAPLES-DR is graded on the Canadian teleophthalmology scale, under which most of its images are R1 (mild) although they were drawn from MESSIDOR grades 1--3, so \rn{clin.maples.lo.n} of its 60 test images fall in grades 0--1 (median \rn{clin.maples.lo.maupr}) and only \rn{clin.maples.hi.n} in grades 2--4 (\rn{clin.maples.hi.maupr}), too few for a comparison. The grade mix of a test set therefore sets the level of every model's score, and the comparisons of this study hold within a test set and not across sets.

\paragraph{Human reference}

DiaRetDB1 gives the comparison a human reference. Each of its four experts, scored against the consensus of the other three with the same protocol, reached a pixel-level \mAUPR{} of \rn{human.expert.maupr.lo}--\rn{human.expert.maupr.hi} (median \rn{human.expert.maupr.median}) and an image-level AUC of \rn{human.expert.auc.lo}--\rn{human.expert.auc.hi} (Table~\ref{M-tab:independent}, bottom block, and Supplementary Table~13; a four-level confidence map yields a four-point precision--recall curve, so the pixel-level value is a lower bound). The retrained models scored against the four-expert consensus (seed 0) spanned \rn{human.models.maupr.lo}--\rn{human.models.maupr.hi} pixel-level \mAUPR{} (median \rn{human.models.DDR.maupr.median} from DDR, \rn{human.models.IDRiD.maupr.median} from IDRiD) and \rn{human.models.auc.lo}--\rn{human.models.auc.hi} image-level AUC (median \rn{human.models.DDR.auc.median} and \rn{human.models.IDRiD.auc.median}), so on pixel-level agreement with the consensus all but \rn{human.models.below} of the \rn{human.models.pairs} model--source pairs lie within or above the band of the experts, whereas on image-level detection only \rn{human.models.DDR.auc.nreach} (DDR-trained) and \rn{human.models.IDRiD.auc.nreach} (IDRiD-trained) of \rn{n.ranked} reach the experts' lowest value of \rn{human.expert.auc.lo}; the experts' own microaneurysm agreement was the lowest of the four classes (AUPR \rn{human.expert.MA.lo}--\rn{human.expert.MA.hi} against \rn{human.models.MA.median} for the median model), which says as much about the consensus rule as about the models.

\paragraph{Ensembles and reference model}

Averaging the probability maps of the retrained models, in compositions fixed before any ensemble was scored (Methods), gave the highest scores of the study on every test set from DDR and on all but one from IDRiD (Table~\ref{M-tab:independent}, lower block; seed-0 members). From DDR, the mean of all 18 members reached 0.505 on the DDR test set and the mean of the five best by validation \mAUPR{} 0.511, against 0.469 for the best member; the same two ensembles reached 0.677 and 0.680 on IDRiD (best member 0.645), 0.637 and 0.652 on the standard-field TJDR set (0.624) and 0.496 and 0.506 on MAPLES-DR (0.466). On the coarsely and region-labelled sets the gain was smaller (0.298 and 0.290 on FGADR and Retinal-Lesions for the five best members against 0.290 and 0.278 for the best member; 0.335 on DiaRetDB1 against 0.325). The gain over the best member was thus 0.02--0.06 \mAUPR{} on the fine-grained test sets and $-$0.01 to 0.02 on the others, and the four adapted foundation models averaged by themselves matched or exceeded the ensemble of all models (0.509 on DDR, 0.676 on IDRiD, 0.657 on the standard-field TJDR set). On FGADR and Retinal-Lesions the two compositions that contain HACDR-Net were computed without it (its maps for those two sets were not available) and are marked in the table. Ensembles are a reference for users and do not enter the rank tests above.

A reference model that follows these findings, RetiSeg-Ref, averages the two main-table models with the highest DDR validation score (RETFound with UPerNet and SAMed), each trained once on the three fine-grained sources merged (599 images; Methods), and will be released with the benchmark. Trained on the merged sources, the two members reached \rn{n8.samed.DDR} (SAMed) and \rn{n8.retfound.DDR} (RETFound with UPerNet) on the DDR test set and their mean \rn{n8.ens.DDR}, against \rn{n8.best.single.ddr.seed0} for the best single DDR-trained model (\rn{n8.best.single.ddr.seed0.model}, seed 0) and \rn{ensall.DDR} for the mean of all 18 DDR-trained members. On the five test sets outside its training sources the averaged model reached \rn{n8.ens.MAPLES-DR} on MAPLES-DR, \rn{n8.ens.FGADR} on FGADR, \rn{n8.ens.Retinal-Lesions} on Retinal-Lesions, \rn{n8.ens.DiaRetDB1} on DiaRetDB1 and \rn{n8.ens.TJDR-uwf} on the ultra-widefield images at the field-of-view scale, within \rn{n8.vs.all.absmax} of the 18-member DDR-trained ensemble on each of the four standard-scale sets (above it on \rn{n8.vs.all.above}, below on \rn{n8.vs.all.below}; \rn{ensall.FGADR}, \rn{ensall.Retinal-Lesions}, \rn{ensall.DiaRetDB1} and \rn{ensall.MAPLES-DR}) with two members in place of 18; on the test sets of the two other merged sources, which are in-source for it, it reached \rn{n8.ens.IDRiD} on IDRiD and \rn{n8.ens.TJDR-std} on the standard-field TJDR set, above every DDR-trained model and ensemble there but below the IDRiD-trained ensemble on IDRiD (0.764), whose members saw each IDRiD image about four times as often under the IDRiD budget. Its scores are a starting point for users of the benchmark and are not ranked (Supplementary Table~11).

\paragraph{Computational cost}

Because the training amount is fixed in pixels, the cost of a model under the unified protocol reflects its throughput and its patch size: one seed on IDRiD took from \rn{cost.min} GPU hours (HRDecoder and U-Net) to \rn{cost.max} (\rn{cost.max.model}), and a DDR seed \rn{cost.ddr.ratio} times as long (median ratio), the slowest being the models that process small patches with sliding windows (FCT, \rn{cost.ddr.max} hours on DDR; SAM-Adapter 26, HACDR-Net 28 and CLC-Net 39 hours). nnU-Net follows its own schedule of 1,000 epochs rather than the pixel budget (7 hours per source, 11--13 for the residual-encoder plan). The two models that led the tables were not the costliest to train: SAMed took 2.1 and 7.4 hours and RETFound with UPerNet 6.1 and 24 hours per seed. In the inference audit (\rn{cost.infer.n} of the \rn{n.main} models; Methods), one forward pass takes \rn{cost.flops.min} GFLOPs (\rn{cost.flops.min.model}) to \rn{cost.flops.max} (\rn{cost.flops.max.model}, whose 224-pixel windows tile the canvas), \rn{cost.infer.min}--\rn{cost.infer.max}\,s per image on an A100 and \rn{cost.mem.min}--\rn{cost.mem.max}\,GB of GPU memory. Here the two leading models differ by an order of magnitude (SAMed \rn{cost.infer.samed}\,s, RETFound with UPerNet \rn{cost.infer.retfound}\,s per image, against \rn{cost.infer.hrnet}\,s for HRNet-W48 and \rn{cost.infer.unet}\,s for U-Net), so the general-purpose high-resolution network gives the same accuracy at a twentieth of the inference time of the ViT-L model. Accuracy was not related to training cost, inference cost or parameter count (Extended Data Fig.~\ref{M-edfig:cost} and Supplementary Table~8).
